\documentclass[12pt]{article}

% Packages
\usepackage[
  paperwidth=11in,
  paperheight=8.5in,
  margin=0.5in
]{geometry}
\usepackage{amsmath, amssymb}
\usepackage{booktabs}
\usepackage{tabularx}

\newcommand{\tacred}{FS-TACRED}
\newcommand{\qwensmall}{Qwen3-4B}
\newcommand{\gemmasmall}{Gemma3-4B}

\begin{document}

% \begin{table}[htbp]
%     \centering
%     \caption{Experiment on \tacred\ using \qwensmall\ and \gemmasmall\ where we optimize the Relation Inference prompt on 1-shot baseline. 
%     The initial prompt and optimized prompts has no additional examples of LLM-generated.}
%     \label{tab:inital_experiment}
% \setlength{\tabcolsep}{.061in}
% \begin{tabular}{l
%                 r@{ \scriptsize$\pm$ }l r@{ \scriptsize$\pm$ }l r@{ \scriptsize$\pm$ }l
%                 r@{ \scriptsize$\pm$ }l r@{ \scriptsize$\pm$ }l r@{ \scriptsize$\pm$ }l}
% \toprule
%   Prompt Optimization Techniques & \multicolumn{6}{c}{\qwensmall} & \multicolumn{6}{c}{\gemmasmall} \\ 
%   \cmidrule(lr){2-7} \cmidrule(lr){8-13}
%   & \multicolumn{2}{c}{P} & \multicolumn{2}{c}{R} & \multicolumn{2}{c}{F1} &
%   \multicolumn{2}{c}{P} & \multicolumn{2}{c}{R} & \multicolumn{2}{c}{F1} \\ \midrule

% Baseline (without additional examples) \\
% Optimized using Feedback from mistakes \\
% Optimized using Feedback from corrects \\
% Optimized using Feedback from mistakes \& corrects \\
%   \bottomrule

% \end{tabular}
% \end{table}

\section{Initial Method and Experiment Details}

We consider a simplified version of our multi-agent framework to establish a strong baseline.
The full system consists of three LLM agents: 
\begin{enumerate}
    \item Example Generator
    \item Relation Inference Agent
    \item Feedback Agent
\end{enumerate}
In the current setting, we evolve only the prompt of the Relation Inference Agent (Agent~2) using feedback from the Feedback Agent (Agent~3).
The Example Generator (Agent~1) is not used.
No additional in-context examples are introduced.
All experiments are therefore conducted under a \textbf{1-shot baseline} configuration.

We update the inference prompt using structured feedback generated by a separate feedback prompt with three variants:
\begin{itemize}
    \item \textbf{Corrects},
    \item \textbf{Mistakes}, and
    \item \textbf{Corrects + Mistakes}.
\end{itemize}


Let $D_{feedback}$ denote the feedback (training) set.
Let $D_{validation}$ denote the validation set.
Let $D_{test}$ denote the test set.
We initialize the system as $S_0$ with a set of prompts
\(
P = \{P_{inference}, P_{feedback}, P_{mutation}\}.
\)
Only the inference prompt $P_{inference}$ is subject to evolution.
The initial population is defined as $Q = \{S_0\}$.


At each iteration, we perform the following steps for $\texttt{MAX\_ITERATION}=20$.

\begin{enumerate}
    \item Sample a system $S_x$ from the population $Q$ based on its validation performance (see Section~\ref{subsec:sampling_system}).
    \item Sample $f$ instances from $D_{feedback}$, with $f=3$ by default.
    \item Run relation inference on the sampled instances using the inference prompt $P_{inference}$ of $S_x$.
    \item Collect structured feedback on the model outputs using a fixed evaluation prompt $P_{evaluation}$.
    \item Use $P_{mutation}$ to update $P_{inference}$ based on the structured feedback, yielding a new system $S_x'$.
    \item Evaluate $S_x'$ on the validation set $D_{validation}$.
    \item Add $S_x'$ to the population $Q$.
\end{enumerate}

After all iterations, we select the system in $Q$ with the highest validation performance.
The selected system is then evaluated on the test set $D_{test}$.

\paragraph{Models and Dataset:}
We conduct experiments using \textbf{Qwen3-4B} and \textbf{Gemma3-4B} as the underlying language models.
All experiments are performed on the \textbf{FS-TACRED} dataset.

\paragraph{Data Format:}
The FS-TACRED validation and test sets follow an episodic few-shot format with a support set and a query.
The validation episodes are constructed from the original validation split.
The feedback set (using the training split of FS-TACRED) is not episodic; we sample two instances, one serving as a 1-shot support example that defines the relation and the other serving as the query.

\subsection{System Sampling}
\label{subsec:sampling_system}

Let the population be $Q = \{S_1, S_2, \dots, S_{|Q|}\}$, where each system $S_i$ has a validation score $v_i$.
To balance exploration and exploitation, we sample a system from $Q$ according to a softmax distribution over validation scores:

\[
p(S_i) = \frac{\exp\left(v_i / \tau\right)}{\sum_{j=1}^{|Q|} \exp\left(v_j / \tau\right)},
\]

where $\tau > 0$ is a temperature parameter controlling exploration; we initialize $\tau = 1$.

% \bibliographystyle{plain}  
% \bibliography{custom}    

\end{document}
